% Seção 1 — Introdução. Não depende de resultados nem de decisões do grupo.
% Para colar no template LaTeX do professor; citações em referencias.bib.

\section{Introdução}

A classificação de imagens consiste em atribuir a uma imagem um rótulo entre um
conjunto finito de categorias. Embora seja uma tarefa trivial para humanos, ela
é difícil de resolver com regras explícitas: um mesmo objeto aparece sob
diferentes pontos de vista, escalas, iluminações, oclusões e fundos, e
categorias distintas podem compartilhar formas e texturas. Abordagens clássicas
dependiam de descritores projetados manualmente seguidos de um classificador,
de modo que a qualidade do sistema ficava limitada pela qualidade das
características escolhidas pelo projetista.

As redes neurais convolucionais (CNNs) mudaram esse cenário ao aprender as
características diretamente dos dados~\cite{lecun1998gradient}. Três
propriedades as tornam adequadas a imagens: a conectividade local, que explora
a correlação entre pixels vizinhos; o compartilhamento de pesos, que reduz o
número de parâmetros e torna a detecção de um padrão independente de sua
posição; e o empilhamento de camadas com \textit{pooling}, que constrói uma
hierarquia de representações, de bordas e texturas nas primeiras camadas a
partes de objetos nas mais profundas. O resultado da AlexNet na competição
ImageNet de 2012~\cite{krizhevsky2012imagenet} consolidou as CNNs como a
abordagem dominante, e arquiteturas posteriores, como a
VGG~\cite{simonyan2015very} e a ResNet~\cite{he2016deep}, mostraram que
profundidade, uso de \textit{kernels} pequenos e conexões residuais têm papel
decisivo no desempenho. Componentes como \textit{Batch
Normalization}~\cite{ioffe2015batch} e \textit{Dropout}~\cite{srivastava2014dropout}
tornaram o treinamento de redes profundas mais estável e menos sujeito a
\textit{overfitting}.

Este trabalho utiliza o CIFAR-10~\cite{krizhevsky2009learning}, um conjunto de
60.000 imagens coloridas de $32 \times 32$ pixels divididas igualmente em dez
classes (avião, automóvel, pássaro, gato, cervo, cachorro, sapo, cavalo, navio
e caminhão), com 50.000 imagens de treinamento e 10.000 de teste. A baixa
resolução e a semelhança visual entre algumas classes, como gato e cachorro ou
automóvel e caminhão, tornam o problema desafiador, enquanto o tamanho moderado
do conjunto permite treinar e comparar várias configurações de rede com
recursos computacionais limitados. Essas características fazem do CIFAR-10 um
\textit{benchmark} adequado para estudar experimentalmente o efeito de
decisões de projeto de uma CNN.

Treinar uma CNN profunda do zero exige muitos dados rotulados e tempo de
processamento. O \textit{transfer learning}~\cite{pan2010survey} contorna essa
limitação reaproveitando uma rede pré-treinada em um conjunto grande e
genérico, tipicamente o ImageNet~\cite{deng2009imagenet}, com cerca de
1,2 milhão de imagens em 1.000 classes. Como as primeiras camadas aprendem
características gerais e as últimas, características mais específicas da tarefa
original~\cite{yosinski2014transferable}, é possível manter parte da rede
congelada e adaptar apenas as camadas finais ou, alternativamente, ajustar a
rede inteira (\textit{fine-tuning}). A escolha entre essas estratégias envolve
um compromisso entre desempenho, custo computacional e risco de
\textit{overfitting}, e a diferença de domínio entre ImageNet e CIFAR-10
(em especial a resolução) torna essa escolha uma questão experimental.

\subsection{Objetivos}

O objetivo deste trabalho é estudar experimentalmente o uso de CNNs para a
classificação do CIFAR-10, em duas frentes: o projeto de uma arquitetura
própria e o uso de redes pré-treinadas por meio de \textit{transfer learning}.
Especificamente, busca-se:

\begin{itemize}
  \item desenvolver uma CNN própria (\textit{baseline}) e avaliar o impacto de
        modificações arquiteturais controladas, alterando uma característica
        por vez;
  \item analisar a relação entre a complexidade das redes (número de
        parâmetros, tempo de treinamento e de inferência) e seu desempenho;
  \item aplicar \textit{transfer learning} com arquiteturas da literatura
        pré-treinadas no ImageNet, comparando diferentes estratégias de
        congelamento e descongelamento de camadas;
  \item comparar as melhores configurações de cada abordagem no conjunto de
        teste, analisando o desempenho por classe e os principais erros de
        classificação.
\end{itemize}

As principais questões investigadas são:
\begin{enumerate}
  \item Quais alterações na arquitetura da CNN própria produzem maior impacto
        no desempenho e no \textit{overfitting}?
  \item Aumentar a capacidade da rede melhora o desempenho, e a que custo
        computacional?
  \item Os modelos pré-treinados apresentam vantagem em relação à CNN própria?
  \item Como o congelamento ou o descongelamento de camadas afeta o desempenho
        do \textit{transfer learning}?
  \item Quais classes são mais difíceis de classificar e quais são mais
        frequentemente confundidas entre si?
\end{enumerate}

% TODO GRUPO: parágrafo final com a organização do texto (ajustar se o
% template já tiver as seções com outros nomes).
O restante do texto está organizado da seguinte forma: a
Seção~\ref{sec:metodos} descreve os materiais e métodos; a
Seção~\ref{sec:resultados} apresenta e discute os resultados; a
Seção~\ref{sec:conclusao} apresenta as conclusões.
